\chapter{Implementation and Snapshots}

\section{Proposed Workflow}
The operational workflow of OmniFAQ AI is engineered as follows:
\begin{enumerate}
    \item \textbf{Search Execution:} The candidate enters a technical topic into the search bar. The input is debounced for 400ms before sending an asynchronous request to \texttt{api.php}.
    \item \textbf{Weighted Matching:} The backend executes a weighted keyword scoring algorithm across question titles, keyword tags, and model answers.
    \item \textbf{Result Resolution:}
    \begin{itemize}
        \item \textit{If matches exist:} The questions are displayed with category badges, difficulty ratings, and view counts.
        \item \textit{If 0 matches exist:} The query is logged as unresolved in the telemetry table, and the user is offered a ``Generate Live AI Questions'' prompt.
    \end{itemize}
    \item \textbf{AI Synthesis and Database Ingestion:} Upon trigger, the backend calls Google Gemini API, validates the returned structured JSON array, permanently inserts the newly synthesized questions into the database, and marks the query as resolved.
    \item \textbf{Active Practice:} The user types a practice response. The client-side evaluator compares the input against domain keyword tokens and provides instant percentage coverage feedback.
\end{enumerate}

\section{Directory Structure}
The project maintains a clean, decoupled architecture:
\begin{verbatim}
omnifaq_ai/
|-- api.php               # Central REST API Router & Controllers
|-- db.php                # Dual-Driver Database Connection & Auto-Seeder
|-- index.html            # Main Single Page Application (SPA)
|-- style.css             # Glassmorphism Design System & Styles
|-- schema.sql            # MySQL Schema & Initial Seed Data
|-- Dockerfile            # Production Container Configuration
|-- render.yaml           # Infrastructure as Code Blueprint
|-- .env                  # Local Environment Variables (Git-Ignored)
|-- report/               # LaTeX Academic Project Report
\end{verbatim}

\section{User Interface}
The application interface is structured as a glassmorphic dashboard featuring:
\begin{itemize}
    \item A top navigation bar with brand logo and authentication state display.
    \item A centralized search box with dynamic suggestion tags.
    \item A responsive question card grid with animated answer reveal accordions.
    \item An interactive practice area with real-time keyword coverage scoring.
    \item An administrative analytics modal showcasing resolution rates and top unresolved queries.
\end{itemize}

\section{Code Implementation}

\subsection{Backend Database Connection (PHP)}
The following snippet demonstrates how \texttt{db.php} implements the dual-driver connection with automatic SQLite fallback and auto-seeding:

\begin{lstlisting}[language=PHP, caption=Dual-Driver Database Connection Logic (db.php)]
function getDbConnection() {
    static $pdo = null;
    if ($pdo !== null) return $pdo;

    // 1. Attempt MySQL connection if host is configured
    if (!empty(DB_HOST)) {
        try {
            $dsn = "mysql:host=" . DB_HOST . ";port=" . DB_PORT . ";dbname=" . DB_NAME . ";charset=utf8mb4";
            $pdo = new PDO($dsn, DB_USER, DB_PASS, [
                PDO::ATTR_ERRMODE => PDO::ERRMODE_EXCEPTION,
                PDO::ATTR_DEFAULT_FETCH_MODE => PDO::FETCH_ASSOC,
                PDO::ATTR_TIMEOUT => 2
            ]);
            return $pdo;
        } catch (Exception $e) { /* Fall through to SQLite */ }
    }

    // 2. Zero-config fallback to embedded SQLite
    $dbFile = __DIR__ . '/database.sqlite';
    $isNew = !file_exists($dbFile) || filesize($dbFile) === 0;
    $pdo = new PDO("sqlite:" . $dbFile, null, null, [
        PDO::ATTR_ERRMODE => PDO::ERRMODE_EXCEPTION,
        PDO::ATTR_DEFAULT_FETCH_MODE => PDO::FETCH_ASSOC
    ]);
    if ($isNew) initSqliteDatabase($pdo);
    return $pdo;
}
\end{lstlisting}

\subsection{Frontend API Communication (JavaScript)}
This snippet illustrates how the client dispatches debounced asynchronous queries to the backend:

\begin{lstlisting}[language=JavaScript, caption=Debounced Search & Asynchronous API Request (index.html)]
function executeSearch(query) {
    $('#faq-skeleton').removeClass('hidden');
    $.ajax({
        url: 'api.php',
        type: 'POST',
        contentType: 'application/json',
        data: JSON.stringify({
            action: 'get_faqs',
            search: query,
            category: currentCategory
        }),
        dataType: 'json',
        success: function(res) {
            $('#faq-skeleton').addClass('hidden');
            if (res.success && res.data.length > 0) {
                renderFaqList(res.data);
            } else {
                renderEmptyStateWithAiPrompt(query);
            }
        }
    });
}
\end{lstlisting}
